%% kapa: exam=1 round=23 year=2012 period=2 form=A subject=accounting subject_ko=회계학 seq=33 number=33
%% source: Q-Net 감정평가사 기출문제 (cst00309) artlSeq=5052580 file=raw/1008/5052580/2012년 제23회 감정평가사 제1차시험 기출문제(2교시 A형).pdf
%% conversion: kordoc 4.24.51 -> tools/gen1.py (auto); review=proofread-sonnet; flags=none
\begin{question}{33}
(주)한국은 단일제품을 대량으로 생산하고 있다. 직접재료는 공정 초기에 모두 투입되고 가공원가는 공정 중에 균등하게 발생한다. 생산 중에는 공손이 발생하는 데 품질검사를 통과한 수량의 10\%에 해당하는 공손수량은 정상공손으로 간주한다. 공손여부는 공정의 50\% 완성시의 검사시점에서 파악된다. 다음 자료를 이용하여 계산한 비정상공손수량은? (단, 물량흐름은 선입선출법을 가정한다.)
\begin{framedtext}
ㆍ 기초재공품수량 2,000단위(완성도 30\%) \par
ㆍ 당기착수량 8,000단위 \par
ㆍ 당기완성량 7,200단위 \par
ㆍ 기말재공품수량 1,500단위(완성도 60\%) \par
\end{framedtext}
\begin{choices}
  \item 380단위
  \item 430단위
  \item 450단위
  \item 520단위
  \item 540단위
\end{choices}
\end{question}
\answer{2}
